\section{Pekerjaan yang Diberikan oleh Perusahaan}
\subsection{Latihan Membuat \textit{Widget} menggunakan Prinsip SDUI}

Tugas pertama diberikan oleh mentor sebagai latihan pengembangan
\textit{widget} berbasis JSON, dikerjakan pada 28--30 Januari 2026 sebagai
bagian dari minggu pertama pelaksanaan magang.\@Proses pengerjaan dilakukan
secara \textit{local} dengan membuat halaman baru khusus untuk latihan (tidak
\textit{push} ke \textit{repository} agar kode tetap bersih). Tersedia
dokumentasi khusus \textit{Front-End} yang dipelajari untuk memahami struktur
dan implementasi pembuatan \textit{widget}.

\begin{lstlisting}[caption=Contoh JSON untuk Membuat \textit{Widget TabView}, captionpos=b, xleftmargin=2cm, xrightmargin=2cm]
{
  "section_class": 3, 
  "section_orientation": 1,
  "section_main_alignment": 3,
  "section_cross_alignment": 3,
  "section_extra" : {
     "with_navigation_button": true,
     "type" : 1,     
  },
  "value": [
     <List of SECTIONS>    
  ]
};
\end{lstlisting}

Pada tahap awal terdapat kesulitan dalam memahami struktur tersebut, namun
pemahaman meningkat seiring waktu. Melalui proses tersebut, berhasil
dikembangkan beberapa jenis \textit{widget} seperti \textit{Text, TextField,
    Image, Radio, Button, Checkbox, TabView, Dropdown,} dan komponen lainnya.

\begin{figure}[H]
    \centering
    \includegraphics[width=12cm]{3.4-latihan.png}
    \caption{Layar Latihan yang Berhasil Dibuat}
\end{figure}

\subsection{Migrasi Beberapa \textit{Widget} ke Repository FE Terbaru}

Setelah memahami penggunaan JSON untuk membangun UI dengan Flutter, diberikan
tugas \#11858 untuk migrasi beberapa \textit{widget} dari \textit{repository}
sebelumnya (Flutter v3.16.9) ke repository terbaru (Flutter v3.38.5). Selain
itu, diterapkan struktur \textit{Model-View-Controller} (MVC) untuk
meningkatkan keterbacaan (readability) serta skalabilitas (scalability) kode
pada \textit{widget} yang dikembangkan. Task ini dikerjakan pada 28 Januari --
2 Februari 2026 dengan durasi pengerjaan 4 hari dan bobot 5 poin KPI.

Migrasi dilakukan terhadap berbagai \textit{widget}, antara lain
\textit{AntrianTV, TextField, AddImageCamera, Map, Card, CodeSnippet, Calendar,
    PieChart, ProgressChart, ProgressBar, BarChart}, dan komponen lainnya.

\begin{figure}[H]
    \centering
    \includegraphics[width=0.475\textwidth]{3.4-migrasi-1.png}
    \hspace{\fill}
    \includegraphics[width=0.475\textwidth]{3.4-migrasi-2.png}
    \caption{Perbandingan Isi dari Misc Widgets pada HealthyOne v1 dan v2}
\end{figure}

\subsection{\textit{Widget Improvements and Bug Fix} di Flutter}

Sebagian besar pekerjaan di sisi Flutter berupa perbaikan dan peningkatan
fungsionalitas \textit{widget} yang sudah ada. Setiap pekerjaan disertai
informasi ID \textit{task} pada OpenProject beserta durasi pengerjaan dan bobot
poin KPI yang diperoleh. Pengelompokan pada subbab ini disusun berdasarkan tema
pekerjaan (Flutter), sehingga urutannya tidak selalu mencerminkan urutan
kronologis tanggal pengerjaan.

\begin{enumerate}
    \item \textit{MultipleChoice/Checkbox} bisa diklik di area teks (\#10074
          --- 1 hari, 1 poin) \\ Sebelumnya, \textit{widget
              MultipleChoice} hanya bisa dipilih dengan mengklik kotaknya saja. Berbeda
          dengan \textit{SingleChoice} yang area teksnya sudah bisa diklik. Perbaikan
          dilakukan dengan menambahkan komponen \textit{InkWell} di area teks agar
          perilakunya konsisten dengan \textit{SingleChoice}.
          \begin{figure}[H]
              \centering
              \includegraphics[width=8cm]{3.4-checkbox.png}
              \caption{\textit{Widget MultipleChoice}}
          \end{figure}

    \item \textit{Accordion} bisa diimplementasikan di \textit{TemplateForm} (\#10120 --- 3 hari, 8 poin) \\
          \textit{Widget Accordion} sebenarnya sudah berjalan di beberapa layar, namun belum bisa digunakan pada layar yang menggunakan \textit{TemplateForm}. Solusinya adalah menambahkan \textit{parser} khusus untuk \textit{Accordion} beserta modelnya di dalam \texttt{template\_form}.dart.
          \begin{figure}[H]
              \centering
              \includegraphics[width=8cm]{3.4-accordion-closed.png}
              \caption{\textit{Widget Accordion} dalam \textit{TemplateForm} dalam kondisi tertutup}
          \end{figure}
          \begin{figure}[H]
              \centering
              \includegraphics[width=8cm]{3.4-accordion-opened.png}
              \caption{\textit{Widget Accordion} dalam \textit{TemplateForm} dalam kondisi terbuka}
          \end{figure}

    \item \textit{MultiDropdown} bisa \textit{flexible} (\textit{searchable}) (\#10319 --- 3 hari, 3 poin) \\
          Tugas ini awalnya ditugaskan ke \textit{Front-End Supervisor}, namun kemudian dialihkan. Modifikasi dilakukan pada mekanisme \textit{text handling} agar \textit{FlexibleDropdown} bisa menerima lebih dari satu pilihan sekaligus dan memperbarui \textit{state}-nya dengan benar.
          \begin{figure}[H]
              \centering
              \includegraphics[width=8cm]{3.4-multidropdown-flexible.png}
              \caption{\textit{Widget Multidropdown} yang \textit{Flexible}}
          \end{figure}

    \item \textit{Line spacing popup AI Scribe} diperbesar 1.5x (\#10735 --- 3 hari, 1 poin) \\
          Juga dialihkan dari \textit{supervisor}. Perubahannya cukup mudah, dengan mengubah \textit{textStyle} dengan ukuran 1.5x pada teks hasil \textit{output} di \textit{popup AI Scribe}.
          \begin{figure}[H]
              \centering
              \includegraphics[width=8cm]{3.4-aiscribe.png}
              \caption{\textit{Popup AI Scribe}}
          \end{figure}

    \item \textit{Bug} teks terpotong di layar Antrean Layanan (HealthyOne) dan Master Distributor (Canvasser) (\#11050 \& \#11059 --- masing-masing 1 hari, 1 poin) \\
          Dua \textit{bug} ini saling berkaitan. Nama distributor dan teks \textit{Subjective}/Keluhan pada layar Antrean Layanan terpotong karena kurangnya ruang vertikal pada \textit{TabView}. Keduanya diselesaikan dengan menambahkan \textit{padding top} pada komponen tersebut.
          \begin{figure}[H]
              \centering
              \includegraphics[width=8cm]{3.4-textfield-terpotong.png}
              \caption{\textit{Widget TextField} yang sudah ditambahkan \textit{padding TabView}}
          \end{figure}

    \item \textit{Font size} tabel bisa dikustomisasi (\#11091 --- 8 poin) \\
          Ini adalah tugas dengan cakupan paling luas di antara semua pekerjaan Flutter. Permintaannya adalah agar ukuran \textit{font} di dalam tabel, termasuk \textit{editable table} dan seluruh \textit{widget} di dalamnya bisa disesuaikan antara dua ukuran: kecil (ukuran \textit{default}) dan besar. Penyelesaiannya dimulai dengan membuat \texttt{TableFontUtils.dart} sebagai pusat \textit{utils} pengontrol ukuran \textit{font}, yang kemudian diintegrasikan satu per satu ke seluruh widget di dalam tabel: \textit{TextField, IconButton, RaisedButton, SingleDropdown, MultiDropdown, QuantityPicker,} hingga \textit{SingleChoice}. Total pengerjaan memakan waktu delapan hari (4--13 Februari 2026), dan menjadi dasar bagi sejumlah \textit{task} lanjutan (\#12809, \#13563, \#13597, \#14127, \#14263) yang dijabarkan pada item-item berikutnya.
          \begin{figure}[H]
              \centering
              \includegraphics[width=8cm]{3.4-fontsizetable-small.png}
              \caption{\textit{Widget ListBased} dengan \textit{font size} terkecil}
          \end{figure}

          \begin{figure}[H]
              \centering
              \includegraphics[width=8cm]{3.4-fontsizetable-large.png}
              \caption{\textit{Widget ListBased} dengan \textit{font size} terbesar}
          \end{figure}

    \item Halaman lebih panjang agar tidak tertutup FAB (\#11708 --- 2 hari, 3 poin) \\
          Beberapa konten di halaman tertutup oleh \textit{Floating Action Button} (FAB).
          Solusinya adalah menambahkan detektor pada \texttt{section\_parser.dart} yang
          secara otomatis menambahkan \textit{extra bottom padding} jika halaman tersebut
          memiliki FAB Button.
          \begin{figure}[H]
              \centering
              \includegraphics[width=8cm]{3.4-fab-bottompadding.png}
              \caption{Layar yang memiliki \textit{FAB} dengan \textit{extra bottom padding}}
          \end{figure}

    \item \textit{Bug dropdown} tidak bisa dipilih setelah \textit{search} di Canvasser (\#11797 --- 1 hari, 3 poin) \\
          Setelah \textit{user} melakukan \textit{search} di \textit{dropdown}, opsi tidak bisa dipilih. Penyebabnya adalah \textit{indexing} yang tidak terlindungi sebelum digunakan, sehingga bisa \textit{out of bound}. Perbaikan dilakukan dengan menambahkan validasi \textit{index} di \texttt{table.dart}.
          \begin{figure}[H]
              \centering
              \includegraphics[width=8cm]{3.4-bugdropdown-cannotselect.png}
              \caption{\textit{Dropdown} yang tidak bisa dipilih}
          \end{figure}

    \item \textit{Button} Tolak Outlet tidak rapi di Canvasser (\#12071 --- 3 hari, poin belum tercatat) \\
          Dari tiga tombol yang ditampilkan dalam satu \textit{card}, satu di antaranya tidak sejajar dengan dua lainnya. Perbaikan dilakukan melalui modifikasi \texttt{raised\_button.dart} dan \texttt{card\_generator\_ext.dart} untuk merapikan \textit{alignment} tombol.
          \begin{figure}[H]
              \centering
              \includegraphics[width=8cm]{3.4-button-alignment.png}
              \caption{Tombol-tombol yang berantakan}
          \end{figure}

    \item \textit{Dropdown} di \textit{editable table} tidak langsung \textit{update height} (\#12108 --- 2 hari, 2 poin) \\
          Ketika \textit{user} mengetik di \textit{dropdown} dan hasil pencarian dari \textit{back-end} sudah muncul, tinggi \textit{dropdown} tidak ikut menyesuaikan dan tetap pendek seperti saat data kosong. Perbaikan dilakukan pada \texttt{single\_dropdown.dart} agar \textit{state height} diperbarui secara instan begitu data berhasil diambil, tanpa perlu menutup dan membuka dropdown kembali.

    \item \textit{Icon Button} terlalu kecil (\#12160 --- 1 hari, 1 poin) \\
          Tugas ini dibuat secara mandiri setelah ditemukan masalah saat mengerjakan tugas lain. Widget \textit{IconButton} tidak memiliki ukuran \textit{default}, sehingga ukurannya mengikuti ukuran \textit{child}-nya dan terlihat terlalu kecil. Solusinya adalah menambahkan nilai \textit{default} 48px untuk \textit{height} dan \textit{width} di \texttt{icon\_button.dart}.

    \item \textit{Bug overflow} saat data tidak ditemukan di \textit{dropdown} (\#12284 --- 2 hari, 1 poin) \\
          Saat pencarian di \textit{dropdown} tidak menghasilkan data, \textit{widget} mengalami \textit{box overflow}. Penyebabnya adalah penggunaan komponen \textit{Flexible} di dalam \texttt{single\_dropdown.dart} yang tidak kompatibel dengan konteks \textit{layout}-nya. Solusinya adalah menghapus komponen \textit{Flexible} tersebut.
          \begin{figure}[H]
              \centering
              \includegraphics[width=8cm]{3.4-erroroverflow-dropdown.png}
              \caption{\textit{Dropdown} yang \textit{error overflow}}
          \end{figure}

    \item Beberapa \textit{widget} tidak mengikuti perubahan \textit{font size} tabel
          (\#12809 --- 2 hari, 2 poin) \\ Sebagai kelanjutan dari Task \#11091, ditemukan
          bahwa beberapa \textit{widget} di dalam \textit{editable table} masih belum
          mengikuti perubahan ukuran \textit{font}, khususnya \textit{TextField} yang
          tidak bisa di-\textit{edit}. Perbaikan dilakukan pada \texttt{table.dart,
              blank\_widget.dart,} dan \texttt{text.dart} dengan menambahkan parameter
          \textit{isForEditTable} agar \textit{font} bisa mendeteksi ukuran yang
          diperlukan.
          \begin{figure}[H]
              \centering
              \includegraphics[width=8cm]{3.4-widgetfont-tertinggal-1.png}
              \caption{\textit{Text} yang belum mengikuti \textit{update font size}}
          \end{figure}

    \item \textit{Prefix} `Rp' tidak mengikuti perubahan \textit{font size} (\#12928 --- 2 hari, 1 poin) \\
          \textit{Prefix} `Rp' pada input nominal memiliki \textit{styling} tersendiri yang terpisah dari teks pada umumnya, sehingga tidak ter-\textit{update} saat ukuran \textit{font} tabel diubah. Penyelesaiannya adalah menambahkan \textit{dynamic styling} pada \textit{prefix} agar mengikuti \textit{state font size} yang aktif.
          \begin{figure}[H]
              \centering
              \includegraphics[width=8cm]{3.4-widgetfont-tertinggal-2.png}
              \caption{\textit{Prefix} `Rp' yang belum mengikuti \textit{update font size}}
          \end{figure}

    \item \textit{DatePicker} tidak bisa digunakan di dalam tabel (\#12934 --- 2 hari, 3 poin) \\
          \textit{DatePicker} pada layar Penerimaan Stok tidak berfungsi saat digunakan di dalam tabel. Penyebabnya adalah \textit{null handling} di \textit{widget DateField} tidak diimplementasikan dengan baik untuk konteks \textit{editable table}. Perbaikan dilakukan dengan menambahkan \textit{ternary operator} untuk mengonversi nilai kosong menjadi \textit{null} sebelum diproses.
          \begin{figure}[H]
              \centering
              \includegraphics[width=8cm]{3.4-datepicker-bug.png}
              \caption{\textit{Widget DatePicker} yang tidak bisa diklik}
          \end{figure}

    \item \textit{Delete button} tidak mengikuti perubahan \textit{font size} (\#13137 --- 2 hari, 1 poin) \\
          Tombol \textit{delete} di dalam tabel tidak menyesuaikan ukurannya saat \textit{font size} diubah karena \textit{padding} dan \textit{height}-nya tidak memperbarui \textit{state} dengan baik. Solusinya adalah membungkus tombol dengan \textit{ValueListenableBuilder} agar selalu sinkron dengan \textit{state} terkini.
          \begin{figure}[H]
              \centering
              \includegraphics[width=8cm]{3.4-widgetfont-tertinggal-3.png}
              \caption{\textit{Widget DeleteButton} yang belum mengikuti \textit{update font size}}
          \end{figure}

    \item \textit{MultiDropdown} mentok saat di perangkat \textit{mobile} di Canvasser (\#13228 --- 2 hari, 2 poin) \\
          Beberapa \textit{dropdown} di Canvasser tidak bisa di-\textit{scroll} sampai bawah ketika dibuka di perangkat \textit{mobile}. Setelah dianalisis, ditemukan adanya beberapa \textit{scrollbar} yang bertumpuk di dalam \textit{widget MultiDropdown}, sehingga saling menghalangi. Solusinya adalah menghapus \textit{SingleChildScrollView} dan \textit{Column} yang bermasalah, menggantinya dengan \textit{Container} biasa karena \textit{ListView.builder} di dalamnya sudah memiliki \textit{scrollbar} sendiri.

    \item \textit{Widget ImageField} (\textit{upload} gambar) bisa mengambil foto dari kamera (\#13240 --- 4 hari, 5 poin) \\
          Sebelumnya, \textit{widget} ini hanya mendukung \textit{upload} dari \textit{file}. \textit{Improvement} ini menambahkan opsi untuk mengambil foto langsung dari kamera perangkat. Implementasinya dengan membuat \textit{widget} baru \textit{`CapturePhotoScreen'} yang mengakses kamera (\textit{webcam}, kamera depan, maupun belakang), kemudian mengintegrasikannya ke dalam \textit{widget ImageField} sehingga \textit{user} bisa memilih antara \textit{upload} dari \textit{file} atau ambil foto langsung.
          \begin{figure}[H]
              \centering
              \includegraphics[width=8cm]{3.4-capturephotoscreen.png}
              \caption{\textit{Widget CapturePhotoScreen} yang baru dibuat}
          \end{figure}

          \begin{figure}[H]
              \centering
              \includegraphics[width=8cm]{3.4-newimagefield.png}
              \caption{\textit{Widget ImageField} dengan fitur Ambil Foto}
          \end{figure}

    \item Bug fix urutan kolom tanggal (\#13508 --- 3 hari, 3 poin) \\ Kolom bertipe
          tanggal pada tabel tidak bisa diurutkan secara ascending maupun descending.
          Ketika tombol sort ditekankan, data tidak berubah urutan sama sekali.
          Penyebabnya adalah logika perbandingan pada fungsi sort belum menangani tipe
          data \textit{date}, \textit{date-range}, dan \textit{date-time} --- hanya tipe
          angka dan string yang sudah ditangani.

          Perbaikan dilakukan dengan menambahkan \textit{sort handler} khusus untuk
          ketiga tipe tanggal tersebut, menggunakan perbandingan berbasis
          \textit{DateTime} untuk menghasilkan urutan yang akurat. Selain itu,
          ditambahkan pula \textit{null handling} agar baris dengan nilai tanggal kosong
          tidak menyebabkan error saat proses pengurutan berjalan.
          %   \begin{figure}[H]
          %       \centering
          %       \includegraphics[width=8cm]{3.58-sort-tanggal.png}
          %       \caption{Widget tabel setelah perbaikan pengurutan kolom tanggal}
          %   \end{figure}

    \item Text widget custom margin dan padding dari backend (\#13549 --- 8 poin) \\
          Sebelumnya, konfigurasi \textit{margin} dan \textit{padding} pada widget
          \textit{Text} dan \textit{Section} hanya bisa dilakukan secara statis di sisi
          \textit{frontend}. Tidak ada parameter yang bisa dikirim dari \textit{backend}
          melalui JSON SDUI untuk mengatur jarak antar elemen secara dinamis.

          Improvement ini menambahkan dukungan parameter \textit{margin} dan
          \textit{padding} pada model SDUI dari widget \textit{Text}, \textit{TextField},
          dan \textit{Section}, sehingga tim \textit{backend} bisa mengatur tata letak
          komponen langsung dari konfigurasi JSON tanpa perlu mengubah kode
          \textit{frontend}. Nilai default tetap dipertahankan agar layar yang tidak
          mengirim parameter ini tidak mengalami perubahan tampilan.

          Pengerjaan task ini memakan waktu empat hari (17--22 April 2026) karena cakupan
          widget yang terdampak cukup luas --- setiap tipe widget perlu ditangani secara
          terpisah sesuai dengan struktur model yang dimilikinya. Task ini merupakan
          salah satu task dengan bobot poin tertinggi (8 poin) selama periode magang.
          %   \begin{figure}[H]
          %       \centering
          %       \includegraphics[width=8cm]{3.59-custom-padding.png}
          %       \caption{Contoh tampilan widget Text dengan custom padding dari backend}
          %   \end{figure}

    \item Icon Konsultasi tidak mengikuti ukuran font (\#13563 --- 1 poin) \\ Ikon yang
          mewakili fitur Konsultasi di beberapa layar tidak menyesuaikan ukurannya ketika
          nilai \textit{font size} aktif diubah melalui kontrol \textit{TableFontUtils}.
          Sementara teks dan elemen lain di sekitarnya sudah responsif, ikon tersebut
          tetap berukuran tetap.

          Perbaikan dilakukan dengan menghubungkan ukuran ikon ke \textit{state}
          \textit{TableFontUtils} yang sama dengan komponen lainnya, sehingga ikon ikut
          berubah ukuran secara konsisten ketika nilai \textit{font size} diperbarui.
          Task ini diselesaikan dalam satu hari karena lokasi perbaikan sudah dapat
          diidentifikasi dengan cepat.

    \item Improvement posisi simbol required (*) pada TextField (\#13597 --- 1 poin) \\
          Pada widget \textit{TextField} yang memiliki komponen \textit{upload button},
          simbol asterisk (*) sebagai penanda \textit{field required} tidak tampil di
          posisi yang benar. Asterisk muncul terpisah dari label teks karena komponen
          \textit{Expanded} di sekitar tombol upload menggeser posisinya.

          Perbaikan dilakukan dengan menghapus komponen \textit{Expanded} tersebut,
          sehingga asterisk kembali tampil tepat di sebelah label teks seperti yang
          diharapkan. Task ini diselesaikan dalam satu hari karena perubahannya bersifat
          \textit{local} pada satu komponen.

    \item Sorting kolom tanggal --- perbaikan lanjutan (\#13612 --- 3 hari, 2 poin) \\
          Sebagai kelanjutan dari Task \#13508, ditemukan bahwa beberapa edge case pada
          pengurutan kolom tanggal masih belum tertangani dengan benar. Ketika nilai
          tanggal pada baris tertentu kosong (\textit{null}), proses sort menghasilkan
          error karena tidak ada \textit{null handling} yang memadai di beberapa titik
          logika.

          Perbaikan dilakukan dengan menambahkan validasi \textit{null} sebelum proses
          perbandingan tanggal dijalankan, serta memastikan baris dengan nilai kosong
          selalu ditempatkan di posisi terakhir --- baik saat sort ascending maupun
          descending.

    \item Alert widget --- perbaikan posisi tampilan (\#13635 --- 3 hari, 1 poin) \\
          Widget \textit{alert} atau \textit{snackbar} yang muncul sebagai notifikasi
          sering kali tertutup oleh elemen lain di bagian bawah layar, seperti
          \textit{navigation bar} atau tombol aksi. Posisi kemunculannya terlalu rendah
          sehingga tidak selalu terlihat oleh pengguna.

          Perbaikan dilakukan dengan mengangkat posisi tampilan \textit{snackbar} agar
          selalu muncul di atas elemen navigasi bawah. Selain itu, ditambahkan offset
          vertikal yang memperhitungkan tinggi komponen di bawahnya secara dinamis.
          %   \begin{figure}[H]
          %       \centering
          %       \includegraphics[width=8cm]{3.60-alert-position.png}
          %       \caption{Posisi alert/snackbar setelah perbaikan}
          %   \end{figure}

    \item Fix spacing TextField (\#13701 --- 1 hari, 1 poin) \\ Beberapa widget
          \textit{TextField} di berbagai layar memiliki jarak (\textit{spacing}) antar
          elemen yang tidak konsisten --- ada yang terlalu rapat dan ada yang terlalu
          renggang dibandingkan desain. Penyebabnya adalah nilai \textit{padding} yang
          berbeda-beda di setiap instansi widget karena belum ada standar yang diterapkan
          secara terpusat.

          Perbaikan dilakukan dengan menyeragamkan nilai \textit{spacing} pada widget
          \textit{TextField} menggunakan konstanta terpusat, sehingga semua instansi
          memiliki jarak yang konsisten tanpa perlu mengubah setiap layar secara
          individual.

    \item Fitur sort tabel hanya berfungsi pada data angka dan tanggal (\#13816 --- 1
          hari, 2 poin) \\ Ditemukan adanya \textit{typo} pada logika perbandingan data
          di dalam \texttt{table.dart} yang menyebabkan fungsi sort tidak bekerja untuk
          tipe data string. Variabel yang dibandingkan menggunakan nama field yang salah
          sehingga hasil perbandingan selalu bernilai sama, membuat urutan data tidak
          berubah.

          Perbaikan dilakukan dengan memperbaiki nama variabel yang salah ketik tersebut.
          Setelah perbaikan, fungsi sort kini berfungsi dengan benar untuk semua tipe
          data yang didukung.

    \item Adjustment font size tabel --- widget tambahan (\#14127 --- poin belum
          tercatat) \\ Sebagai kelanjutan dari sistem \textit{TableFontUtils} yang
          dibangun pada Task \#11091, ditemukan beberapa widget di dalam \textit{editable
              table} yang masih belum mengikuti perubahan ukuran font secara konsisten.
          Widget-widget ini tidak terdaftar dalam sistem pengontrol ukuran font yang
          sudah ada.

          Perbaikan dilakukan dengan mengidentifikasi semua widget yang belum
          terintegrasi, kemudian menghubungkannya ke \textit{TableFontUtils} menggunakan
          pola yang sama dengan widget-widget yang sudah diimplementasikan sebelumnya.
          Pengerjaan memakan waktu empat hari (20--25 Mei 2026) karena jumlah widget yang
          perlu diperiksa cukup banyak.

    \item Handle ukuran field input saat edit tabel (\#14263 --- 3 poin) \\ Ditemukan
          ketidakkonsistenan tinggi widget dan ukuran ikon pada beberapa komponen Flutter
          di dalam \textit{editable table} yang tidak tertangani oleh sistem
          \textit{TableFontUtils} sebelumnya. Komponen yang terdampak antara lain
          \textit{TextForm, TimePicker, SingleDropdown,} dan \textit{MultiDropdown} ---
          yang masing-masing memiliki logika \textit{sizing} tersendiri yang tidak
          terhubung ke sistem \textit{TableFontUtils}.

          Penyelesaiannya dimulai dengan menganalisis pola ketidakkonsistenan pada setiap
          komponen secara terpisah, kemudian mengimplementasikan penyesuaian
          \textit{height, iconSize,} dan ukuran teks agar konsisten dengan
          \textit{TableFontUtils} yang aktif. Pengerjaan memakan waktu empat hari (18--21
          Mei 2026) karena setiap komponen perlu ditangani secara independen sesuai
          strukturnya masing-masing.
          %   \begin{figure}[H]
          %       \centering
          %       \includegraphics[width=8cm]{3.61-textfield-size.png}
          %       \caption{Widget TextField setelah perbaikan ukuran pada editable table}
          %   \end{figure}

    \item Data hilang setelah filter Group By di Pembuatan Konten AI (\#13670 --- 4 hari,
          1 poin) \\ Pada layar Pembuatan Konten AI, data yang tampil di
          \textit{ListBased} menghilang setelah pengguna mengaktifkan filter
          \textit{Group By}. Penyebabnya adalah fitur \textit{group by} belum didukung
          oleh tipe tampilan \textit{card-vertical} pada \textit{ListBased}, sehingga
          ketika \textit{group by} diaktifkan dalam mode tampilan tersebut, komponen
          mencoba melakukan render dengan struktur data yang tidak kompatibel.

          Perbaikan dilakukan dengan menonaktifkan opsi \textit{group by} secara otomatis
          ketika \textit{ListBased} berada dalam mode tampilan \textit{card-vertical},
          sehingga pengguna tidak dapat mengaktifkan kombinasi yang belum didukung.

    \item Pesan validasi field birthdate Data Pasien tertimpa (\#14699 --- 1 hari, 1
          poin) \\ Pada layar Data Pasien, pesan validasi untuk field tanggal lahir
          (\textit{birthdate}) tertimpa oleh widget lain yang berada di sekitarnya karena
          komponen pesan error tidak memiliki batasan lebar yang memadai. Selain itu,
          simbol \textit{required} (*) pada \textit{text-form} ditampilkan dengan warna
          yang tidak sesuai dalam beberapa kondisi.

          Perbaikan dilakukan dengan membungkus komponen pesan error menggunakan
          \textit{Flexible} agar ukurannya menyesuaikan ruang yang tersedia dan tidak
          menimpa elemen lain, serta memperbaiki logika pewarnaan simbol
          \textit{required} di widget \textit{text-form}.

    \item Dropdown tidak lengkap menampilkan data di Riwayat Pengobatan (\#14858 --- 2
          hari, poin belum tercatat) \\ Widget \textit{search-chip} pada layar Riwayat
          Pengobatan tidak menampilkan semua data yang tersedia. Ketika daftar item cukup
          panjang, bagian bawah \textit{dropdown} terpotong dan tidak bisa
          di-\textit{scroll}. Penyebabnya adalah properti \textit{prototypeItem} yang
          mengunci tinggi komponen ke nilai tetap berdasarkan satu item contoh, sehingga
          komponen tidak bisa melakukan ekspansi tinggi secara dinamis.

          Perbaikan dilakukan dengan menghapus properti \textit{prototypeItem}, sehingga
          tinggi komponen \textit{search-chip} bisa menyesuaikan diri secara dinamis
          sesuai jumlah item yang ditampilkan.

    \item Diagnosis Sekunder hilang setelah interaksi AI Scribe (\#14374 --- poin belum
          tercatat) \\ Pada layar konsultasi, data Diagnosis Sekunder yang sudah diisi
          oleh dokter menghilang setelah pengguna berinteraksi dengan fitur AI Scribe.
          Penyebabnya adalah \textit{items provider} pada widget \textit{search\_chip}
          tidak menangani kondisi di mana AI Scribe melakukan \textit{update state} pada
          komponen secara bersamaan --- ketika AI Scribe memperbarui \textit{state},
          nilai \textit{items provider} direset ke kondisi awal sehingga data yang sudah
          ada ikut terhapus.

          Perbaikan dilakukan dengan menambahkan mekanisme penjagaan pada \textit{items
              provider} agar tidak direset saat AI Scribe melakukan \textit{update},
          melainkan hanya diperbarui ketika ada perubahan eksplisit dari pengguna.
          Pengerjaan memakan waktu tujuh hari (27 Mei -- 4 Juni 2026) karena reproduksi
          bug ini memerlukan interaksi spesifik antara dua fitur yang berbeda.

\end{enumerate}

Poin KPI yang diperoleh dari seluruh \textit{task} pada subbab ini, ditambah
dengan pengerjaan di sisi React pada subbab berikutnya, direkap secara bulanan
pada Subbab 3.3.5.

\subsection{Pembuatan Layar-layar HealthyOne di React}
Berbeda dengan pekerjaan di Flutter yang didominasi oleh perbaikan dan
peningkatan \textit{widget}, pengembangan di React lebih berfokus pada
pembuatan layar baru untuk aplikasi HealthyOne. Pengembangan layar Data Pasien
dan Item Cabang dilakukan pada \textit{repository}
\texttt{healthyone-medic-fe}, sedangkan modul-modul manajemen stok (Penerimaan
Stok, Stok Opname, Perpindahan Stok, Riwayat Stok, Pengemasan dan Pengeluaran,
Riwayat Belanja, Dashboard Inventaris, dan DataView) serta Marketplace B2B
dikembangkan pada \textit{repository} \texttt{healthy-one-stock-management-fe}.

\begin{enumerate}
    \item Data Pasien (\#12200, \#12202, \#12203 --- 16--20 Februari 2026, total 10 poin)
          \begin{figure}[H]
              \centering
              \includegraphics[width=10cm]{screen-healthyone-flutter-datapasien.png}
              \caption{Tampilan layar Data Pasien di Flutter}%
          \end{figure}

          Modul ini terdiri dari empat layar utama: layar daftar pasien yang menampilkan
          data dalam format tabel, layar \textit{View} dengan dua \textit{tab} (Biodata
          Pasien dan Rekam Medis), layar \textit{Edit} dengan form yang sudah terisi data
          pasien, serta layar Tambah dengan form kosong untuk \textit{input} data baru.
          Layar ini juga dilengkapi fitur \textit{split screen}, sehingga form tambah,
          edit, atau lihat bisa ditampilkan berdampingan dengan tabel tanpa perlu
          berpindah halaman.

          \begin{figure}[H]
              \centering
              \includegraphics[width=10cm]{screen-healthyone-react-datapasien.png}
              \caption{Tampilan layar Data Pasien di React}%
          \end{figure}

    \item Item Cabang (\#12127, \#12116 --- awal Maret 2026, total 6 poin)
          \begin{figure}[H]
              \centering
              \includegraphics[width=10cm]{3.4-itemcabang-flutter.png}
              \caption{Tampilan layar Data Pasien di Flutter}%
          \end{figure}

          Terdapat empat tab utama dengan isi masing-masing sebagai berikut:
          \begin{enumerate}
              \item Tab Produk menampilkan enam subtab: Semua, Obat Resep, Produk Kesehatan, Produk
                    BMHP, Produk Lainnya, dan Obat Racik. Khusus untuk lima subtab pertama (kecuali
                    Obat Racik), masing-masing memiliki lima mode tampilan: Daftar Produk,
                    Rekomendasi Belanja (Defecta), Stok Opname, Penjualan Terbesar, dan Riwayat
                    Kedaluwarsa. Subtab Obat Racik memiliki logika tersendiri karena tidak
                    menggunakan mekanisme filtering yang sama. Setiap subtab juga mendukung operasi
                    Tambah, View, dan Edit produk. Konfigurasi tampilan Stok Opname dan Rekomendasi
                    Belanja pada tab ini masing-masing dikerjakan melalui Task \#12343 dan \#12344
                    (1 hari, 2 poin masing-masing, 11 Maret 2026).
                    \begin{figure}[H]
                        \centering
                        \includegraphics[width=10cm]{3.4-itemcabang-react-produk.png}
                        \caption{Tampilan layar Data Pasien bagian Produk di React}%
                    \end{figure}
              \item Tab Layanan menampilkan empat subtab: Konsultasi, Gigi, Lainnya, dan Belum
                    Aktif. Masing-masing subtab menampilkan daftar layanan yang tersedia di unit
                    pelayanan, lengkap dengan operasi Tambah, Edit, dan View.
                    \begin{figure}[H]
                        \centering
                        \includegraphics[width=10cm]{3.4-itemcabang-react-layanan.png}
                        \caption{Tampilan layar Data Pasien bagian Layanan di React}%
                    \end{figure}
              \item Tab Tindakan menampilkan daftar tindakan yang tersedia di unit pelayanan,
                    dengan fitur Tambah Tindakan.
                    \begin{figure}[H]
                        \centering
                        \includegraphics[width=10cm]{3.4-itemcabang-react-tindakan.png}
                        \caption{Tampilan layar Data Pasien bagian Tindakan di React}%
                    \end{figure}
              \item Tab Paket menampilkan daftar paket (bundle) yang tersedia, dengan dukungan
                    operasi Tambah, Edit, dan Lihat Paket masing-masing dalam bentuk form
                    tersendiri.
                    \begin{figure}[H]
                        \centering
                        \includegraphics[width=10cm]{3.4-itemcabang-react-paket.png}
                        \caption{Tampilan layar Data Pasien bagian Paket di React}%
                    \end{figure}
          \end{enumerate}

    \item Penerimaan Stok
          \begin{figure}[H]
              \centering
              \includegraphics[width=10cm]{3.4-penerimaanstok-flutter.png}
              \caption{Tampilan layar Penerimaan Stok di Flutter}%
          \end{figure}

          Modul ini terdiri dari tiga layar: layar utama Penerimaan Stok, layar View
          Penerimaan, dan layar Tambah Penerimaan. Selain ketiga layar tersebut, dibuat
          pula komponen Jenis Penerimaan Stok yang digunakan sebagai bagian dari alur
          input data penerimaan. Pengembangan modul ini tercakup dalam Task \#12576
          (modifikasi FE manajemen inventaris, 10 hari, 8 poin, 18--31 Maret 2026)
          bersama modul Stok Opname, Perpindahan Stok, Riwayat Stok, dan Riwayat Belanja
          pada item-item berikutnya.
          \begin{figure}[H]
              \centering
              \includegraphics[width=10cm]{3.4-penerimaanstok-react.png}
              \caption{Tampilan layar Penerimaan Stok di React}%
          \end{figure}

    \item Stok Opname (\#12117 --- 6 hari, 8 poin)
          \begin{figure}[H]
              \centering
              \includegraphics[width=10cm]{3.4-stokopname-flutter.png}
              \caption{Tampilan layar Stok Opname di Flutter}%
          \end{figure}

          Layar Stok Opname pada dasarnya merupakan layar Riwayat Stok yang dirender
          ulang dengan parameter query \texttt{isForStokOpname=true}, disertai tambahan
          tombol Buat Stok Opname di bagian antarmukanya.
          \begin{figure}[H]
              \centering
              \includegraphics[width=10cm]{3.4-stokopname-react.png}
              \caption{Tampilan layar Stok Opname di React}%
          \end{figure}

          Pengembangan modul ini dilanjutkan dengan proses \textit{alignment} ke
          \textit{backend} Toraja (Task \#13410, 8 hari, 8 poin) yang menghubungkan
          antarmuka React dengan API. Selama fase pengujian di \textit{environment
              staging}, dilakukan serangkaian penyempurnaan dan perbaikan \textit{bug} yang
          meliputi:
          \begin{itemize}
              \item \textbf{Perbaikan Sistem Seleksi dan Pagination (\#14201 \& \#14281)}: Memperbaiki bug di mana seleksi produk hilang saat berpindah halaman (\textit{pagination}). Perbaikan dilakukan dengan menyimpan baris terpilih ke dalam \textit{state} \texttt{allLoadedRows} lintas halaman. Masalah konflik \textit{state} seleksi produk ketika membuka dua halaman Stok Opname secara bersamaan juga diselesaikan dengan mengisolasi \textit{state} \texttt{selectedIds} pada masing-masing komponen.
              \item \textbf{Validasi Form dan Pesan Error (\#14282 \& \#14128)}: Menambahkan validasi pop-up untuk mencegah pengguna menekan tombol \textit{submit} tanpa memilih produk terlebih dahulu, serta mengganti pesan error bawaan menjadi komponen \textit{Alert} terstandarisasi.
              \item \textbf{Konsistensi Kontrol Input dan Dropdown (\#14131, \#14129, \& \#14486)}: Menonaktifkan dropdown Unit Pelayanan karena nilainya sudah ditentukan otomatis dari sesi pengguna. Menghilangkan \textit{overlap} scrollbar pada \textit{QuantityPicker} dengan mengubah pengaturan \texttt{zIndex}, serta memperbaiki bug konversi kuantitas unit terkecil pada komponen \textit{pilih-nomor-lot} di dalam \texttt{useEffect}.
              \item \textbf{Penyelarasan Tipe Data dan Sorting (\#14133, \#14125, \#14176, \& \#14130)}: Menyeragamkan format tanggal menjadi standar Indonesia (\texttt{dd MMM yyyy}), memperbaiki ketidaksesuaian penamaan kolom pada payload API, menonaktifkan kontrol \textit{sorting widget} untuk kolom yang belum mendukung pengurutan berbasis \textit{backend}, serta menambahkan penanganan khusus untuk tombol Stok Opname pada kategori produk lain-lain.
          \end{itemize}

    \item Perpindahan Stok
          \begin{figure}[H]
              \centering
              \includegraphics[width=10cm]{3.4-perpindahanstok-flutter.png}
              \caption{Tampilan layar Perpindahan Stok di Flutter}%
          \end{figure}

          Serupa dengan Stok Opname, layar ini juga memanfaatkan kembali komponen Riwayat
          Stok dengan query yang sesuai, ditambah tombol Ajukan Perpindahan Stok dan
          komponen notifikasi sebagai pembeda fungsionalnya. Pengembangan layar ini juga
          tercakup dalam Task \#12576.
          \begin{figure}[H]
              \centering
              \includegraphics[width=10cm]{3.4-perpindahanstok-react.png}
              \caption{Tampilan layar Perpindahan Stok di React}%
          \end{figure}

    \item Riwayat Stok
          \begin{figure}[H]
              \centering
              \includegraphics[width=10cm]{3.4-riwayatstok-flutter.png}
              \caption{Tampilan layar Riwayat Stok di Flutter}%
          \end{figure}

          Terdiri dari dua layar: layar utama Riwayat Stok dan layar Lihat Detail
          Penerimaan Stok. Layar ini juga menjadi basis yang digunakan kembali oleh modul
          Stok Opname dan Perpindahan Stok. Pengembangan layar ini juga tercakup dalam
          Task \#12576.
          \begin{figure}[H]
              \centering
              \includegraphics[width=10cm]{3.4-riwayatstok-react.png}
              \caption{Tampilan layar Riwayat Stok di React}%
          \end{figure}

    \item Pengemasan dan Pengeluaran (\#12130 --- 1 hari, 2 poin)
          \begin{figure}[H]
              \centering
              \includegraphics[width=10cm]{3.4-pengemasanpengeluaran-flutter.png}
              \caption{Tampilan layar Pengemasan dan Pengeluaran di Flutter}%
          \end{figure}

          Modul ini mengelola alur pemrosesan pesanan dari awal hingga selesai. Terdapat
          enam tab utama yang merepresentasikan status pesanan: Pesanan Baru, Disiapkan,
          Siap Dikirim, Dalam Pengiriman, Dibatalkan, dan Pesanan Selesai. Masing-masing
          tab memiliki subtab berdasarkan jenis pesanan (Online dan/atau Offline) sesuai
          dengan alur distribusi yang relevan untuk setiap status.
          \begin{figure}[H]
              \centering
              \includegraphics[width=10cm]{3.4-pengemasanpengeluaran-react.png}
              \caption{Tampilan layar Pengemasan dan Pengeluaran di React}%
          \end{figure}

    \item Riwayat Belanja
          \begin{figure}[H]
              \centering
              \includegraphics[width=10cm]{3.4-riwayatbelanja-flutter.png}
              \caption{Tampilan layar Riwayat Belanja di Flutter}%
          \end{figure}

          Modul ini terdiri dari tiga layar: layar utama Riwayat Belanja, layar Tambah
          Pesanan dengan form input, dan layar Lihat Pesanan untuk melihat detail
          transaksi yang sudah ada. Pengembangan layar ini juga tercakup dalam Task
          \#12576.
          \begin{figure}[H]
              \centering
              \includegraphics[width=10cm]{3.4-riwayatbelanja-react.png}
              \caption{Tampilan layar Riwayat Belanja di React}%
          \end{figure}

    \item Riwayat Belanja Marketplace B2B (\#12579 --- 4 hari, 3 poin, 23--26 Maret 2026)
          \begin{figure}[H]
              \centering
              \includegraphics[width=10cm]{3.4-riwayatbelanjab2b-flutter.png}
              \caption{Tampilan layar Riwayat Belanja B2B di Flutter}%
          \end{figure}

          Layar ini menampilkan riwayat transaksi B2B yang dibagi menjadi dua tab: Belum
          Lunas dan Sudah Lunas. Masing-masing tab memiliki layar Detail Belanja untuk
          melihat informasi transaksi secara lengkap.
          \begin{figure}[H]
              \centering
              \includegraphics[width=10cm]{3.4-riwayatbelanjab2b-react.png}
              \caption{Tampilan layar Riwayat Belanja B2B di React}%
          \end{figure}

    \item Marketplace B2B (\#12798 --- 3 hari, 3 poin, 30 Maret -- 1 April 2026)
          \begin{figure}[H]
              \centering
              \includegraphics[width=10cm]{3.4-marketplaceb2b-flutter.png}
              \caption{Tampilan layar Marketplace B2B di Flutter}%
          \end{figure}

          Modul terakhir yang dikerjakan adalah Marketplace B2B, yang terdiri dari lima
          layar: layar utama Marketplace, layar Healthy Coin, layar Produk Kesehatan,
          layar Cart, dan layar Detail Produk. Modul ini juga dilengkapi beberapa fitur
          tambahan, yaitu upload tanda tangan untuk pengajuan File SP, serta pemilihan
          alamat pengiriman, jenis ekspedisi, dan metode pembayaran dalam alur checkout.
          \begin{figure}[H]
              \centering
              \includegraphics[width=10cm]{3.4-marketplaceb2b-react.png}
              \caption{Tampilan layar Marketplace B2B di React}%
          \end{figure}

    \item Sidebar psychotest-management-app (\#13590 --- 4 hari, 5 poin, 20--23 April
          2026) \\ Sebagai \textit{additional task} di luar repositori HealthyOne utama,
          turut dikerjakan pembangunan komponen Sidebar untuk aplikasi
          \textit{psychotest-management-app} --- sebuah aplikasi internal perusahaan yang
          digunakan untuk manajemen tes psikologi calon karyawan. Komponen Sidebar ini
          mencakup navigasi antar halaman, tampilan informasi pengguna yang sedang login,
          dan \textit{routing} ke modul-modul utama aplikasi.

    \item Kolom keterangan pada form Stok Opname (\#12566 --- 24 April 2026) \\ Pada form
          pembuatan Stok Opname, ditambahkan kolom input \textit{keterangan} yang
          memungkinkan pengguna memberikan catatan atau penjelasan tambahan saat membuat
          entri Stok Opname baru. Penambahan ini dilakukan langsung di environment
          \textit{staging} karena bersifat penyesuaian kecil pada form yang sudah ada.

    \item Layar Detail Klinik Canvasser (\#14481 --- 4 hari, 5 poin, 19--22 Mei 2026) \\
          Layar Detail Klinik merupakan layar baru yang dibangun untuk aplikasi Canvasser
          di sisi React. Layar ini menampilkan informasi lengkap mengenai satu klinik
          mitra, termasuk data identitas klinik, daftar produk yang tersedia, daftar
          layanan, dan tombol-tombol aksi yang relevan untuk tenaga lapangan.
          %   \begin{figure}[H]
          %       \centering
          %       \includegraphics[width=10cm]{3.62-detail-klinik.png}
          %       \caption{Tampilan layar Detail Klinik Canvasser di React}
          %   \end{figure}

    \item Dashboard Inventaris --- Laporan Stok Opname dan Laporan Pembelian (\#13286,
          \#13592, \#13999, \#14234 --- total 18 poin, akhir April--Mei 2026) \\
          Dashboard Inventaris merupakan modul baru yang menjadi pusat pelaporan untuk
          pengelolaan inventaris klinik. Modul ini terdiri dari dua sub-modul utama yang
          dapat ditampilkan secara berdampingan melalui panel \textit{resizable split
              screen}:

          \textbf{(a) Laporan Stok Opname} menampilkan data historis seluruh aktivitas Stok Opname yang pernah dilakukan di klinik, dengan dukungan \textit{filtering} berdasarkan rentang tanggal, nama produk, status, dan berbagai parameter lainnya. Data ditampilkan menggunakan komponen \textit{DataView} (lihat item 15) yang mendukung pengurutan, pengelompokan, dan ekspor data.

          \textbf{(b) Laporan Pembelian} menampilkan riwayat transaksi pembelian inventaris klinik beserta statusnya. Layar ini dilengkapi dengan tombol Bayar yang memungkinkan admin klinik memproses pembayaran langsung dari halaman yang sama tanpa harus berpindah halaman.

          Selain dua sub-modul tersebut, Dashboard Inventaris juga mendukung fitur
          \textit{filtering} antarmuka yang fleksibel: pengguna dapat menyaring data
          berdasarkan beberapa kriteria sekaligus, termasuk filter lintas kolom yang
          dikombinasikan dengan \textit{group by} lebih dari satu kolom (Task \#14234, 6
          hari, 11--18 Mei 2026).

          Selama fase pengembangan dan migrasi modul ini, diselesaikan beberapa perbaikan
          fungsionalitas:
          \begin{itemize}
              \item \textbf{Penanganan Navigasi Aksi Pembayaran (\#14627)}: Mengubah perilaku tombol Bayar pada Laporan Pembelian agar pengguna tetap berada di halaman Dashboard Inventaris setelah pembayaran berhasil diproses, alih-alih dialihkan ke halaman baru.
              \item \textbf{Optimasi Validitas Data Tabel (\#14572)}: Mengatasi masalah hilangnya data tabel pada Laporan Pembelian akibat pemanggilan ulang API (\textit{re-fetch}) yang dipicu perubahan \textit{state} internal. Struktur pemeriksaan \texttt{laporanType} diperbaiki untuk menstabilkan data.
              \item \textbf{Perbaikan UI Pasca-Migrasi (\#14646)}: Memperbaiki masalah \textit{overflow wrap} tampilan \textit{CardView} yang terlalu sempit pada layar \textit{mobile}, serta memperbaiki bug pada penghapusan \textit{filter chips}.
          \end{itemize}
          %   \begin{figure}[H]
          %       \centering
          %       \includegraphics[width=10cm]{3.63-stok-opname.png}
          %       \caption{Dashboard Inventaris tampilan Laporan Stok Opname}
          %   \end{figure}
          %   \begin{figure}[H]
          %       \centering
          %       \includegraphics[width=10cm]{3.64-laporan-pembelian.png}
          %       \caption{Dashboard Inventaris tampilan Laporan Pembelian}
          %   \end{figure}

    \item Komponen DataView (\#14855 --- 23 hari, 30 April--1 Juni 2026) \\ DataView
          merupakan komponen generik baru yang dibangun dari awal sebagai pengganti
          \textit{ListBased} yang sudah ada, dengan \textit{library} \textbf{AG Grid}
          sebagai basis implementasinya. Pembangunan komponen ini dipicu oleh
          meningkatnya kebutuhan akan tabel yang mendukung berbagai mode tampilan,
          interaksi tingkat lanjut, dan konfigurasi yang fleksibel dari setiap halaman
          pemakainya. Task ini merupakan \textit{task} dengan durasi pengerjaan terlama
          selama periode magang yang dilaporkan.

          Fitur-fitur utama yang dimiliki \textit{DataView} antara lain:
          \begin{itemize}
              \item \textbf{TableView dan CardView} yang dapat dipilih secara dinamis oleh pengguna.
              \item \textbf{Backend filtering} dengan berbagai tipe: \textit{text search, single select, multi select,} dan \textit{date range}.
              \item \textbf{Backend group by} dengan dukungan pengelompokan berdasarkan lebih dari satu kolom sekaligus.
              \item \textbf{Visible columns}: pengguna dapat menyembunyikan atau menampilkan kolom tertentu sesuai kebutuhan.
              \item \textbf{Selectable rows} untuk keperluan aksi massal.
              \item \textbf{Export CSV dan XLSX} langsung dari antarmuka tabel.
              \item \textbf{Custom render components} yang mengikuti \textit{font size provider} aktif.
              \item \textbf{Auto-formatting kolom} melalui fungsi \texttt{generateColumnDef} yang otomatis mendeteksi tipe data dan menerapkan format yang sesuai (\texttt{Rp} untuk kolom \textit{currency}, format tanggal Indonesia untuk kolom \textit{date}).
          \end{itemize}

          Pembangunan komponen \textit{DataView} mencakup penyelesaian beberapa sub-task
          penting:
          \begin{itemize}
              \item \textbf{Ekspor Data Terintegrasi (\#14696)}: Mengimplementasikan fitur ekspor data ke format CSV dan XLSX langsung dari sisi \textit{client}. Proses ekspor dilengkapi fungsi validasi \texttt{hasActiveFilters} untuk menjamin keselarasan data yang diunduh dengan filter aktif di layar, serta menonaktifkan pewarnaan baris selang-seling saat penulisan file Excel.
              \item \textbf{Pengelompokan Multi-Kolom (\#14697 \& \#14740)}: Mengembangkan fitur pengelompokan data (\textit{group-by}) berbasis beberapa kolom secara bersamaan. Logika pengurutan pada kolom \texttt{jumlah\_data} saat pengelompokan aktif diperbaiki agar tidak menghasilkan kesalahan urutan.
              \item \textbf{Format Data Otomatis (\#14965)}: Menyempurnaan fungsi \texttt{generateColumnDef} dengan deteksi otomatis tipe data untuk memformat angka nominal langsung dengan simbol mata uang rupiah (\texttt{Rp}) serta kolom tanggal dengan format lokal Indonesia.
              \item \textbf{Kompatibilitas Library Komponen (\#14579)}: Memperbaiki masalah \textit{state checkbox} pada \textit{FilterMenu} yang tidak langsung diperbarui ketika diklik akibat perubahan versi pustaka Mantine UI yang digunakan di proyek.
          \end{itemize}

          Setelah komponen selesai dibangun, dilakukan migrasi bertahap dari
          \textit{ListBased} ke \textit{DataView} pada seluruh halaman yang
          menggunakannya, yaitu \textit{laporan-stok-opname, laporan-pembelian,
              dashboard-inventaris,} dan \textit{pilih-nomor-lot}. Proses migrasi ini
          dikerjakan melalui Task \#14586 (4 hari, 5 poin, 25--28 Mei 2026).
          %   \begin{figure}[H]
          %       \centering
          %       \includegraphics[width=10cm]{3.65-dataview-table.png}
          %       \caption{Komponen DataView dalam tampilan TableView}
          %   \end{figure}
          %   \begin{figure}[H]
          %       \centering
          %       \includegraphics[width=10cm]{3.66-dataview-card.png}
          %       \caption{Komponen DataView dalam tampilan CardView dengan filter aktif}
          %   \end{figure}

    \item Landing Page Stock Management (28 Mei -- 8 Juni 2026) \\ Sebelum landing page
          ini dibuat, pengguna yang mengakses aplikasi \textit{Stock Management} langsung
          diarahkan secara otomatis ke halaman tertentu tanpa diberi pilihan. Hal ini
          kurang fleksibel karena berbagai modul (Riwayat Stok, Stok Opname, Dashboard
          Inventaris, dan lainnya) memiliki kebutuhan akses yang berbeda-beda tergantung
          peran pengguna.

          Landing page yang dibuat menampilkan semua modul yang tersedia dalam satu
          halaman beranda, lengkap dengan \textit{routing} ke masing-masing modul.
          Bersamaan dengan pembuatan landing page ini, seluruh file \textit{ListBased}
          lama yang sudah tidak digunakan dihapus dari repositori sebagai bagian dari
          \textit{cleanup} pasca-migrasi ke \textit{DataView}.
          %   \begin{figure}[H]
          %       \centering
          %       \includegraphics[width=10cm]{3.67-landing-page.png}
          %       \caption{Tampilan Landing Page Stock Management}
          %   \end{figure}

\end{enumerate}

\subsection{Alignment Front-End ke Back-End}

Mulai pertengahan April 2026 (ditandai dengan Task \#13410, 20--29 April 2026),
Tim FE ditugaskan untuk menghubungkan layar-layar React yang telah dibuat
dengan API dari Back-End agar data dapat ditampilkan secara dinamis. Seiring
berjalannya pekerjaan ini, peran Front-End berkembang menjadi Full-Stack,
karena pekerjaan tidak lagi hanya berada di sisi React, tetapi juga melibatkan
pembuatan dan penyesuaian \textit{endpoint} di sisi Back-End (Toraja).

Pada tahap awal proses \textit{alignment}, ditemukan kendala dalam menyiapkan
\textit{local development environment} Toraja, terutama terkait ketidaksesuaian
versi Python dan beberapa \textit{dependency} yang dibutuhkan. Kendala ini
diselesaikan dengan bantuan tim Back-End dan menjadi dasar saran perbaikan
dokumentasi \textit{setup environment} yang disampaikan pada Subbab 5.2.

Di sisi React, dikembangkan beberapa \textit{layer} pendukung seperti
\textit{services, hooks, mapper, helper,} dan \textit{utils} untuk menangani
proses \textit{fetch data}, pengolahan \textit{response}, hingga distribusi
data ke komponen. Sementara itu, di sisi Toraja, dibuat dan disesuaikan
berbagai API untuk menampilkan halaman React melalui \textit{iframe}, serta API
yang digunakan oleh React untuk mengambil data dari database.

\begin{enumerate}
    \item \textbf{Riwayat Stok} (Task \#14000 --- 5 hari, 2 poin, 29 April--5 Mei 2026) \\
          Halaman React ditampilkan melalui API \texttt{GetRiwayatStokView} menggunakan \textit{iframe widget}. Setelah halaman dimuat, data riwayat stok diambil melalui API \texttt{GetStockHistoryTableData}, kemudian ditampilkan dalam bentuk tabel pada halaman tersebut.

    \item \textbf{Stok Opname} (Task \#13410 --- 8 hari, 8 poin, 20--29 April 2026) \\
          Prosesnya serupa dengan Riwayat Stok. Halaman React ditampilkan melalui API \texttt{GetPilihNomorLotView} menggunakan \textit{iframe widget}. Untuk pengambilan data, digunakan API \texttt{GetDynamicStockOpnameView} yang dikembangkan oleh Tim Back-End, sehingga data dapat dimuat dan ditampilkan sesuai kebutuhan. Selain itu, dikembangkan pula mekanisme navigasi \textit{split screen} antara React dan Toraja melalui fungsi \texttt{redirectSplitScreen} (Task \#13848, 29 April 2026), yang memungkinkan halaman React mengirim sinyal navigasi ke Toraja melalui \textit{iframe} tanpa menyebabkan reload halaman keseluruhan. Pada saat laporan ini ditulis, Task \#13848 masih berstatus belum selesai.

    \item \textbf{Dashboard Inventaris --- Manajemen Inventaris Medis} (Task \#13999 --- 5 hari, 8 poin, 30 April--6 Mei 2026) \\
          Di antara seluruh pekerjaan \textit{alignment} yang dilakukan, Dashboard Inventaris menjadi yang paling menandai pergeseran peran dari Front-End menjadi Full-Stack secara penuh.

          Berbeda dengan Riwayat Stok dan Stok Opname yang hanya memerlukan penyesuaian
          ringan di sisi Backend, \textit{alignment} Dashboard Inventaris melibatkan
          pembuatan dan modifikasi endpoint di Toraja secara lebih aktif. Di sisi React,
          dikembangkan \textit{layer services, hooks,} dan \textit{mapper} untuk
          menangani komunikasi antara komponen dan API, mulai dari proses \textit{fetch}
          data hingga pengolahan \textit{response} yang diterima. Di sisi Toraja, dibuat
          beberapa endpoint baru --- termasuk untuk \texttt{laporan-stok-opname} dan
          \texttt{laporan-pembelian} --- yang dikonfigurasi untuk menyediakan data sesuai
          kebutuhan tampilan.

          Proses ini menuntut koordinasi dengan tim Backend yang cukup intensif, karena
          beberapa endpoint perlu dibangun dari awal atau dimodifikasi agar sesuai dengan
          kebutuhan data di sisi tampilan. Keterlibatan aktif di sisi Toraja pada fase
          ini menjadi awal dari transisi peran Tim Front-End menuju Full-Stack,
          sebagaimana yang disebutkan dalam rencana pengembangan tim di Subbab 3.2.
          %   \begin{figure}[H]
          %       \centering
          %       \includegraphics[width=10cm]{3.68-dashboard-connected.png}
          %       \caption{Dashboard Inventaris yang sudah terkoneksi dengan Back-End}
          %   \end{figure}

    \item \textbf{Riwayat Belanja dan Details Belanja} (Task \#14857 dan \#15085 --- mulai 3 dan 8 Juni 2026, masih berjalan) \\
          Sebagai kelanjutan dari proses \textit{alignment} yang sedang berjalan, pekerjaan berikutnya adalah menghubungkan halaman Riwayat Belanja dan Details Belanja React dengan endpoint Toraja yang sesuai. Proses ini sudah dimulai pada 3 dan 8 Juni 2026 dan masih berlangsung hingga saat laporan ini ditulis. Pola \textit{alignment} yang diterapkan mengikuti pendekatan yang sama dengan modul-modul sebelumnya: halaman React ditampilkan melalui \textit{iframe} di Toraja, dan data diambil melalui endpoint yang dibuat atau disesuaikan oleh Tim Back-End.

    \item \textbf{Perbaikan Bug dan Integrasi Repository Lintas Fungsi} (Task \#14553 \& \#14271) \\
          Selain menghubungkan modul React, peran \textit{Full-Stack} juga diimplementasikan dalam perbaikan logika bisnis di repositori \textit{back-end} lainnya:
          \begin{itemize}
              \item \textbf{Reparasi Penampilan Data Canvasser (\#14553)}: Melakukan pembenahan pada \textit{repository} \texttt{healthcare-productdetection} di modul \texttt{pssp\_program/request\_pssp/data\_process.py}. Bug berupa warna latar belakang kolom teks yang berubah menjadi abu-abu saat berisi karakter \texttt{'-'} diselesaikan dengan menghapus flag \texttt{is\_currency} pada kolom alasan pengajuan PSSP.
              \item \textbf{Perbaikan Duplikasi Data Supplier (\#14271)}: Menyelesaikan masalah \textit{double entry} saat proses penyimpanan detail data supplier pada dokumen PSSP barang di sistem \textit{back-end} utama.
          \end{itemize}
\end{enumerate}

\subsection{Perkiraan Pekerjaan yang Akan Dilanjutkan}

Ke depan, pekerjaan utama masih berpusat pada migrasi dan \textit{alignment}
HealthyOne serta Canvasser ke React. Beberapa modul yang belum selesai
di-\textit{align} ke Back-End --- termasuk Riwayat Belanja, Details Belanja,
dan Penerimaan Stok --- akan menjadi prioritas berikutnya segera setelah
laporan ini selesai disusun. Proses \textit{alignment} ini dikerjakan secara
\textit{full-stack}, artinya pengerjaan tidak terbatas pada sisi tampilan React
saja, tetapi juga mencakup pembuatan dan penyesuaian endpoint di Toraja.

Selain itu, masih terdapat satu task lama yang belum terselesaikan, yaitu Task
\#13636 (\textit{saat print, jangan diarahkan ke new tab}), yang pengerjaannya
sempat tertunda karena adanya prioritas task lain yang lebih mendesak.

Di sisi Flutter, masih ada potensi perbaikan lanjutan terkait konsistensi
sistem \textit{TableFontUtils} yang mungkin ditemukan seiring dengan pengujian
lebih lanjut oleh tim QA.

Project Manager juga telah menyampaikan rencana pengembangan proyek baru
bernama \textbf{Principal Suite}. Informasi terkait proyek ini masih bersifat
\textit{confidential}, sehingga belum dapat dijelaskan secara rinci dalam
laporan ini. Namun, Tim Front-End direncanakan akan terlibat secara
\textit{full-stack} dalam proyek tersebut dengan \textit{tech stack} yang
berbeda dari proyek-proyek sebelumnya. Pengerjaan proyek ini akan dimulai
setelah proses migrasi sistem yang sedang berjalan sudah cukup stabil.

Sebagai catatan tambahan, pemberian poin KPI untuk sejumlah \textit{task} yang
dijabarkan pada subbab-subbab sebelumnya (\#12071, \#14127, \#14858, dan
\#14374) masih dalam proses validasi oleh mentor pada saat laporan ini disusun,
sehingga belum tercantum pada rekapitulasi poin bulanan di Subbab 3.3.5.
Pembaruan rekapitulasi tersebut akan disampaikan pada laporan magang periode
selanjutnya.